\section{Limits and Stability of the Index under Sharing}
\label{sec:robust}

Sections~\ref{sec:main} and~\ref{sec:comparative} assumed a per-arm planning law. Under
sharing, two results mark where index optimality ends, and a regret bound shows when the index
remains near-optimal. Costs are those of Assumption~\ref{assm:standing}: an arm's next cost may
depend on every arm's history. Define $G_i(S)=G_i^{(\nu_i)}(S)$ by advancing only $i$ from the
current history, with other arms' observations, and their effect on $i$'s costs, fixed along
this hypothetical run; the index policy recomputes these indices after each observation. Shared observations can change untouched arms' laws and costs;
freezing hyperparameters does not restore independence.

\begin{proposition}[No universally optimal own-arm-law policy]\label{prop:no-local}
There are two instances with the same incumbent, costs, available stages and separate laws of
each arm's remaining readings, one independent and one shared, with different uniquely optimal
first actions. Consequently no single rule whose inputs consist only of those quantities is
optimal for all joint planning laws, including randomized rules.
\end{proposition}

The construction couples a noncollectable signal in one arm to another arm's collectable
reward while preserving their separate laws; for one-shot boxes,
\citet{bao2026interdependent} show similarly that fixed box indices cannot be optimal under
interdependence. A positive standalone run value still certifies that stopping is suboptimal,
and under a forward-closure condition the index stopping region is exactly optimal
(Proposition~\ref{prop:stop-any}, Appendix~\ref{app:sharing}).

\paragraph{Collection reshapes acquisition under a shared state.}
Sharing also interacts with designated collection: the same joint law of readings can
yield different acquisition orders under different collection rights.

\begin{proposition}[Collection-induced acquisition switching]\label{prop:switching}
There is a two-arm family with a common binary latent state of prior $p$, a nondegenerate
closed interval of terminal costs, and thresholds $0<\tau_1<\tau_2(c)<\tau_3(c)<1$ depending
continuously on the cost $c$ in that interval, with the following property. If arm $B$'s
intermediate reading is collectable, whether or not $A$'s is, the uniquely optimal first
action advances $A$ for $p\in(0,\tau_1)\cup(\tau_2(c),\tau_3(c))$ and advances $B$ for
$p\in(\tau_1,\tau_2(c))\cup(\tau_3(c),1)$. If it is not collectable, whether or not $A$'s is,
advancing $A$ is the uniquely optimal first action for every $p\in(0,1)$.
\end{proposition}

In the construction, keeping $B$'s larger intermediate score adds an option to avoid a terminal
inspection, which competes with $A$'s terminal reward and information.

A prior-dependent change of first inspection already occurs in
one-shot interdependent search \citep[Example~1]{bao2026interdependent}. The family,
its thresholds, a fixed-cost first-action table and the proof are in
Appendix~\ref{app:switching}; the proof is an exact
computer-assisted Bellman certificate over the field of rational functions of $c$, uniform
over the cost interval.

\paragraph{Stability under a coherent shared law.}
A quantitative guarantee follows by specializing approximate-planning bounds
\citep[Theorem~9 and Remark~11]{subramanian2022ais} to the independent value representation.
Restrict further: all readings are coordinates of complete arm paths with one fixed joint law
$\mathsf P$, and predictions are its conditional laws; dependence within and across paths is
arbitrary. For arm $j$ at $S$, let $c_{j,S}(s)$ be the cost of its next stage once only $j$ has
advanced from $S$ to own state $s$, as in the runs defining $G_j(S)$. Let $\mathsf P_{j,S}$ be
the posterior marginal law of arm $j$'s remaining path at $S$, $Z_j^S$ its covered call under
that law and the costs $c_{j,S}$, and $v_{j,S}(r)=\E_{\mathsf P_{j,S}}\max\{r,Z_j^S\}$ its local
keep-or-run value. For the actual
successor $S'=F(S,i,Y)$, the discrepancy
\[
 \varepsilon_i(S)=\E\!\left[
   \sum_{j\ne i}\sup_{r\in[q(S'),y_{\max}]}
       |v_{j,S'}(r)-v_{j,S}(r)|\,\middle|\,S,i\right]
\]
measures how much an observation changes untouched arms' local values, through their laws
and their costs; each supremum is the
stop-loss metric between the two capped variables clipped at the new incumbent
\citep[Section~5, Example~C]{muller1997ipm}.

\begin{proposition}[Index regret under sharing]\label{prop:sharing-robustness}
Under Assumption~\ref{assm:standing} and the coherent-law restriction above, let $H=L(S_0)$. If $\varepsilon_i(S)\le\bar\varepsilon$ almost surely
along every admissible policy from $S_0$, then the policy $\pi_{\rm IX}$,
recomputed from the true posterior marginal laws after every observation, satisfies
\[
 0\le V^\star(S_0)-V_{\pi_{\rm IX}}(S_0)\le 2H\bar\varepsilon.
\]
\end{proposition}

The bound covers premature stops and incorrect advances and needs no separation between
indices. With independent paths and own-history costs $\varepsilon_i=0$, recovering
Theorem~\ref{thm:index}; a cost that moves with another arm's reading enters $\varepsilon_i$
even when the paths are independent. The
envelope $\bar\varepsilon$ measures how one reading moves the untouched arms' local
keep-or-run values above the new incumbent, not correlation: a shared nuisance that shifts
remaining-path laws without moving those values leaves it at zero
(Appendix~\ref{app:sharing-controls}). It must hold at every feasible future history,
including states at which the index would already have stopped, so small sharing at the root
is insufficient: in Appendix~\ref{app:sharing-robustness} arm paths are pairwise independent
at the root and every root envelope is zero, yet the index stops while buying a signal has
positive value. The bound controls neither posterior misspecification nor numerical index
error; Appendix~\ref{app:sharing-robustness} proves it through an explicit signed sharing
correction and gives an envelope through path total variation and cost changes.

\paragraph{Sharing through a learned common component.}
In Bayesian optimization, sharing typically flows through a component common to all
candidates, such as a task-level parameter of a learning-curve model learned from previous
tasks and from the current run. Suppose arm paths are conditionally independent given a latent
$\Theta$ on a standard Borel space, with posterior $\mu_S$ at state $S$, so that
$\mathsf P_{j,S}=\int\mathsf P_j^\vartheta(\cdot\mid s_j)\,\mu_S(d\vartheta)$ for
measurable component path kernels $\mathsf P_j^\vartheta$. Let
\[
 \Lambda_j(S)=\sup_{\vartheta,\vartheta'}\operatorname{TV}\bigl(\mathsf P_j^\vartheta(\cdot\mid s_j),\mathsf P_j^{\vartheta'}(\cdot\mid s_j)\bigr),\qquad
 \eta_i(S)=\E_{\vartheta\sim\mu_S}\operatorname{TV}\bigl(\mathsf P_i^\vartheta(Y\in\cdot\mid s_i),\mathsf P_{i,S}(Y\in\cdot)\bigr),
\]
be the \emph{sensitivity} of arm $j$'s remaining path to the component and the posterior
\emph{dispersion} of arm $i$'s next reading, with
$\operatorname{TV}(\mu,\nu)=\sup_A|\mu(A)-\nu(A)|$, and let
$B_j(S)=y_{\max}-y_{\min}+\nu_j(S)c_{\max}$ if arm $j$ has a collectable stage left and
$B_j(S)=0$ otherwise. Let each cost be a posterior mean
$c_{j,S}(s)=\E[k_j(\Theta,s)\mid S,s]$ of a measurable component cost
$k_j\in[c_{\min},c_{\max}]$, and let $\Lambda^c_j(S)=\sup|k_j(\vartheta,s)-k_j(\vartheta',s)|$
over components and arm $j$'s continuation states, or zero if no collectable stage is left.
Assume the sensitivities are measurable;
this is automatic for a finite component space.

\begin{corollary}[Learned common component]\label{cor:component}
Under Assumption~\ref{assm:standing} and the coherent-law restriction of Proposition~\ref{prop:sharing-robustness}, with arm paths conditionally independent given $\Theta$ and component costs as above, write $\beta_j=B_j\Lambda_j+\nu_j\Lambda^c_j$. Then
\[
 \varepsilon_i(S)\le\eta_i(S)\sum_{j\ne i}\beta_j(S),
 \qquad\text{so}\qquad
 0\le V^\star(S_0)-V_{\pi_{\rm IX}}(S_0)\le 2H\sup_{S,i}\eta_i(S)\sum_{j\ne i}\beta_j(S)
\]
over feasible states and advances. Moreover
$\eta_i(S)\le\E_{\vartheta,\vartheta'\sim\mu_S\otimes\mu_S}
\operatorname{TV}(\mathsf P_i^\vartheta(Y\in\cdot\mid s_i),\mathsf P_i^{\vartheta'}(Y\in\cdot\mid s_i))$,
the expected disagreement of two posterior draws. For binary $\Theta$ with
$p=\mu_S(\Theta=1)$, $\eta_i(S)=2p(1-p)\operatorname{TV}(\mathsf P_i^0(Y\in\cdot\mid s_i),\mathsf P_i^1(Y\in\cdot\mid s_i))$
and $\Lambda_j(S)$ is the separation of the two component path laws.
\end{corollary}

Sensitivity measures how far a candidate's law and costs can move with the component,
dispersion how much one more reading can move the belief about it.
For a finite component space, posterior mass tending to one on one component makes
the local dispersion vanish. With continuous components, posterior concentration needs
additional regularity, such as total-variation continuity of the next-reading kernels;
weak concentration alone does not suffice. Either way, the envelope must be controlled at
rare future histories whose beliefs may be less concentrated. Appendix~\ref{app:sharing-robustness} gives a noiseless
counterexample to automatic dispersion decay.
Appendix~\ref{app:sharing-robustness} gives the proof and Section~\ref{sec:validation} the
checks.
